\section{Coupling constraints: exact output and deferred proofs}\label{app:tu}

This appendix proves the exact-output results of Section~\ref{sec:tu-exact}
and gives the details for Section~\ref{sec:tu-limits}. The notation is that of
Section~\ref{sec:tu-exact}. For $z\in\mathcal Z$ put $\phi_z(x)=F(x,z)$, so that
$\Delta\nabla\phi_z(x)=\hat H_{xx}x+c_z$; for a set $\mathcal J$ of rows,
$\hat A_{\mathcal J}$ consists of the rows of $\hat A$ in $\mathcal J$. If $z$
is integral and $x\in\rho^{-1}\Z^{n_c}$ for a positive integer $\rho$, then
$\Delta\rho^2F(x,z)\in\Z$, since $\Delta$ clears every monomial coefficient.

\subsection{Stationary face polytopes and rational height}\label{app:tu-height}

\begin{definition}[Stationary face polytope]\label{def:tu-statpoly}
For $\tilde v=(\tilde x,z)\in\mathcal S$ let
$\mathcal J_0=\{r:\hat A_r\tilde x=(\hat b_z)_r\}$ be its active rows and
\[
\mathcal P(\tilde v)=\bigl\{x\in\mathcal P_z:\
\hat A_{\mathcal J_0}x=(\hat b_z)_{\mathcal J_0},\ \
\nabla\phi_z(x)\perp\ker\hat A_{\mathcal J_0}\bigr\}.
\]
\end{definition}

\begin{lemma}[Stationary face polytopes]\label{lem:tu-statpoly}
Let $\tilde v=(\tilde x,z)\in\mathcal S$ and $\mathcal J_0$ be as above.
\begin{enumerate}[label=(\alph*)]
\item $\nabla\phi_z(\tilde x)\perp\ker\hat A_{\mathcal J_0}$, and
$w^{\T}\hat H_{xx}w\ge0$ for $w\in\ker\hat A_{\mathcal J_0}$.
\item $\mathcal P(\tilde v)$ is a nonempty polytope and
$\mathcal P(\tilde v)\times\{z\}\subseteq\mathcal S$.
\item Every vertex $x^\circ$ of $\mathcal P(\tilde v)$ lies in
$\rho^{-1}\Z^{n_c}$ for an integer $\rho$ with $\Delta_\eta\mid\rho$ and
$\rho\le\Delta_\eta R_{\mathrm{TU}}$.
\end{enumerate}
\end{lemma}

\begin{proof}
(a) For $w\in\ker\hat A_{\mathcal J_0}$ and small $\lvert t\rvert$,
$\tilde x+tw\in\mathcal P_z$: active rows stay active and inactive rows keep
positive slack. The function
$t\mapsto\phi_z(\tilde x+tw)=\OPT+t\nabla\phi_z(\tilde x)^{\T}w+\frac{t^2}{2}w^{\T}H_{xx}w$
has a minimum at $t=0$, which gives both claims, since
$\hat H_{xx}=\Delta H_{xx}$ with $\Delta>0$.

(b) By (a), $\tilde x\in\mathcal P(\tilde v)$. The condition
$\nabla\phi_z(x)\perp\ker\hat A_{\mathcal J_0}$ reads
$\Xi_0^{\T}(\hat H_{xx}x+c_z)=0$ for a basis matrix $\Xi_0$ of
$\ker\hat A_{\mathcal J_0}$, which is linear, and
$\mathcal P(\tilde v)\subseteq\mathcal P_z$ is bounded. For
$x\in\mathcal P(\tilde v)$ put $w=x-\tilde x\in\ker\hat A_{\mathcal J_0}$. Both
$\nabla\phi_z(x)$ and $\nabla\phi_z(\tilde x)$ are orthogonal to
$\ker\hat A_{\mathcal J_0}$, hence so is their difference $H_{xx}w$, and
\[
\phi_z(x)-\phi_z(\tilde x)=\nabla\phi_z(\tilde x)^{\T}w+\tfrac12w^{\T}H_{xx}w=0 .
\]
Since $(x,z)\in X$, it is a minimizer.

(c) Let $\mathcal J_{x^\circ}\supseteq\mathcal J_0$ be the rows active at
$x^\circ$, let $\mathcal E\subseteq\mathcal J_{x^\circ}$ index linearly
independent rows of $\hat A$ that span the rows of
$\hat A_{\mathcal J_{x^\circ}}$. Then
$\ker\hat A_{\mathcal E}=\ker\hat A_{\mathcal J_{x^\circ}}\subseteq\ker\hat A_{\mathcal J_0}$.

First, $\hat H_{xx}$ is positive definite on $\ker\hat A_{\mathcal E}$. Indeed,
let $w\in\ker\hat A_{\mathcal E}$ with $w^{\T}\hat H_{xx}w=0$. By (a) the
quadratic form of $\hat H_{xx}$ is nonnegative on $\ker\hat A_{\mathcal J_0}$
and vanishes at $w$, so for $y\in\ker\hat A_{\mathcal J_0}$ and all real $t$,
$0\le(w+ty)^{\T}\hat H_{xx}(w+ty)=2t\,y^{\T}\hat H_{xx}w+t^2y^{\T}\hat H_{xx}y$,
whence $\hat H_{xx}w\perp\ker\hat A_{\mathcal J_0}$. Then
$x^\circ\pm tw\in\mathcal P(\tilde v)$ for small $t>0$: rows in
$\mathcal J_{x^\circ}$ stay active, the other rows keep positive slack, and the
stationarity condition is preserved. As $x^\circ$ is a vertex, $w=0$.

Next, $\nabla\phi_z(x^\circ)\perp\ker\hat A_{\mathcal J_0}\supseteq\ker\hat A_{\mathcal E}$,
so $\hat H_{xx}x^\circ+c_z=\hat A_{\mathcal E}^{\T}\pi$ for some vector $\pi$.
Thus $(x^\circ,-\pi)$ solves
\[
\mathsf K\begin{pmatrix}x^\circ\\ -\pi\end{pmatrix}
=\begin{pmatrix}-c_z\\ (\hat b_z)_{\mathcal E}\end{pmatrix},\qquad
\mathsf K=\begin{pmatrix}\hat H_{xx}&\hat A_{\mathcal E}^{\T}\\ \hat A_{\mathcal E}&0\end{pmatrix}.
\]
The matrix $\mathsf K$ is nonsingular: if $\hat H_{xx}w+\hat A_{\mathcal E}^{\T}\pi'=0$
and $\hat A_{\mathcal E}w=0$, then $w^{\T}\hat H_{xx}w=-(\hat A_{\mathcal E}w)^{\T}\pi'=0$,
so $w=0$, and then $\hat A_{\mathcal E}^{\T}\pi'=0$ gives $\pi'=0$ because
$\hat A_{\mathcal E}$ has full row rank. The matrix $\mathsf K$ is integral and
the right-hand side lies in $\Delta_\eta^{-1}\Z^{n_c+\lvert\mathcal E\rvert}$,
so Cramer's rule gives $x^\circ\in\rho^{-1}\Z^{n_c}$ with
$\rho=\Delta_\eta\lvert\det\mathsf K\rvert$, where $\lvert\det\mathsf K\rvert$
is a positive integer. By Hadamard's inequality for rows
\cite{HornJohnson2013},
$\lvert\det\mathsf K\rvert$ is at most the product of the row norms of
$\mathsf K$. A row of
$(\hat H_{xx}\ \hat A_{\mathcal E}^{\T})$ has $n_c$ entries of absolute value at
most $C_{\hat H}$ and $\lvert\mathcal E\rvert\le n_c$ entries in $\{0,\pm1\}$,
so its norm is at most $\sqrt{2n_c}\,C_{\hat H}$; a row of
$(\hat A_{\mathcal E}\ 0)$ is a row of $\hat A$ and has norm at most
$\sqrt{n_c}$. Hence
$\lvert\det\mathsf K\rvert\le(2n_cC_{\hat H}^2)^{n_c/2}n_c^{\lvert\mathcal E\rvert/2}
\le(\sqrt2\,n_cC_{\hat H})^{n_c}\le R_{\mathrm{TU}}$.
\end{proof}

The argument follows Vavasis \cite{Vavasis1990}, with explicit
constants that use only the continuous Hessian block and the $\{0,\pm1\}$
structure of $\hat A$.

\begin{corollary}[Rational height with TU coupling]\label{cor:tu-height}
(a) Some global minimizer has continuous coordinates in $\rho^{-1}\Z$ for an
integer $1\le\rho\le\Delta_\eta R_{\mathrm{TU}}$, and every isolated global
minimizer has this property.
(b) $\OPT=a/W$ in lowest terms with $1\le W\le\Omega_{\mathrm{TU}}$.
\end{corollary}

\begin{proof}
(a) For $\tilde v\in\mathcal S$ the polytope $\mathcal P(\tilde v)$ has a
vertex, which is a minimizer by Lemma~\ref{lem:tu-statpoly}(b) and has the
stated height by Lemma~\ref{lem:tu-statpoly}(c). If $\tilde v$ is isolated in
$\mathcal S$, the convex set $\mathcal P(\tilde v)\times\{z\}\subseteq\mathcal S$
contains no other point near $\tilde v$, so $\mathcal P(\tilde v)=\{\tilde x\}$
and $\tilde x$ is its vertex.
(b) Evaluating $F$ at a minimizer of height $\rho\le\Delta_\eta R_{\mathrm{TU}}$
gives $\Delta\rho^2\OPT\in\Z$, so $W\le\Delta(\Delta_\eta R_{\mathrm{TU}})^2=\Omega_{\mathrm{TU}}$.
\end{proof}

Corollary~\ref{cor:tu-height} is an explicit-constant instance of the
classical fact that quadratic programs over polyhedra have optimal solutions
of polynomial size \cite{Vavasis1990}; for mixed-integer quadratic programs see
\cite{DelPiaDeyMolinaro2017}.

\subsection{Snapping on rows}\label{app:tu-snap}

\begin{lemma}[Snapping recovery on TU faces]\label{lem:tu-snap}
Let $(y,z)\in X$ with $\dist((y,z),\mathcal S)\le\tau_{\mathrm{TU}}/(2\sqrt{n_c})$,
and let $\mathcal J$ and $\mathcal P_z(\mathcal J)$ be as in steps~(i)
and~(ii) of TU-EXACT (Algorithm~\ref{alg:tuexact}). Then $\mathcal P_z(\mathcal J)\ne\emptyset$ and
$\mathcal P_z(\mathcal J)\times\{z\}\subseteq\mathcal S$.
\end{lemma}

\begin{proof}
Write $\tau=\tau_{\mathrm{TU}}$. Let $\tilde v$ be a point of $\mathcal S$
nearest to $(y,z)$; it exists because $\mathcal S$ is compact. Since
$\norm{(y,z)-\tilde v}<1$, the integer parts agree, so $\tilde v=(\tilde x,z)$
and $\norm{y-\tilde x}\le\tau/(2\sqrt{n_c})$. Let $\mathcal J_0$ be the rows
active at $\tilde x$. Every row of $\hat A$ has norm at most $\sqrt{n_c}$, so
for $r\in\mathcal J_0$,
$0\le(\hat b_z)_r-\hat A_ry=\hat A_r(\tilde x-y)\le\tau/2$; hence
$\mathcal J_0\subseteq\mathcal J$. For $r\in\mathcal J\setminus\mathcal J_0$ the
slack at $\tilde x$ is at most $\tau+\tau/2$. On $\mathcal P(\tilde v)$ the
linear function
\[
\psi(x)=\sum_{r\in\mathcal J\setminus\mathcal J_0}\bigl((\hat b_z)_r-\hat A_rx\bigr)
\]
is nonnegative, and
$\psi(\tilde x)\le m'\cdot\frac32\tau=\frac3{8\Delta_\eta R_{\mathrm{TU}}}
<\frac1{\Delta_\eta R_{\mathrm{TU}}}$. Its minimum over the polytope
$\mathcal P(\tilde v)$ is attained at a vertex $x^\circ$. By
Lemma~\ref{lem:tu-statpoly}(c), $x^\circ\in\rho^{-1}\Z^{n_c}$ for an integer
$\rho$ with $\Delta_\eta\mid\rho$ and $\rho\le\Delta_\eta R_{\mathrm{TU}}$; as
$\hat A$ is integral and
$\hat b_z\in\Delta_\eta^{-1}\Z^{m'}\subseteq\rho^{-1}\Z^{m'}$,
$\psi(x^\circ)\in\rho^{-1}\Z$. Since
$0\le\psi(x^\circ)<1/(\Delta_\eta R_{\mathrm{TU}})\le1/\rho$,
$\psi(x^\circ)=0$.

So $x^\circ$ satisfies every row of $\mathcal J\setminus\mathcal J_0$ with
equality, and every row of $\mathcal J_0$ because $x^\circ\in\mathcal P(\tilde v)$.
Moreover $\ker\hat A_{\mathcal J}\subseteq\ker\hat A_{\mathcal J_0}$ and
$\nabla\phi_z(x^\circ)\perp\ker\hat A_{\mathcal J_0}$, so
$\Xi_{\mathcal J}^{\T}(\hat H_{xx}x^\circ+c_z)=0$. Hence
$x^\circ\in\mathcal P_z(\mathcal J)$. For any $x\in\mathcal P_z(\mathcal J)$ let
$w=x-x^\circ\in\ker\hat A_{\mathcal J}$. Both $\nabla\phi_z(x)$ and
$\nabla\phi_z(x^\circ)$ are orthogonal to $\ker\hat A_{\mathcal J}$, hence so is
$H_{xx}w$, and
$\phi_z(x)=\phi_z(x^\circ)+\nabla\phi_z(x^\circ)^{\T}w+\frac12w^{\T}H_{xx}w=\phi_z(x^\circ)=\OPT$,
the last equality by Lemma~\ref{lem:tu-statpoly}(b). Since
$x\in\mathcal P_z$, $(x,z)\in\mathcal S$.
\end{proof}

No nonsingularity, strict complementarity or sign condition on multipliers is
used: the condition $\Xi_{\mathcal J}^{\T}(\hat H_{xx}x+c_z)=0$ allows
multipliers of either sign. A point returned from a distant $y$ can be a
nonglobal stationary point: for $F(x)=x_1-x_1^2$ on
$\{x\in[0,1]^2:x_1=x_2\}$ and $y=(\frac12,\frac12)$, steps~(i) and~(ii)
return $(\frac12,\frac12)$, with value $\frac14>\OPT=0$. Acceptance therefore
always goes through the test in step~(iii).

\subsection{Exact output}\label{app:tu-exact-proof}

\begin{proof}[Proof of Theorem~\ref{thm:tu-exact}]
(a) By Proposition~\ref{prop:tu-sound}(e), $\beta_j\le\OPT$, and
$(\hat x,z)\in X$ because $\hat x\in\mathcal P_z$. By
Corollary~\ref{cor:tu-height}(b), $\OPT$ is a rational with denominator at most
$\Omega_{\mathrm{TU}}$, so Proposition~\ref{prop:accept} with
$\Omega'=\Omega_{\mathrm{TU}}$ shows that a point that passes the test in
step~(iii) is a global minimizer.

(b) By Proposition~\ref{prop:tu-sound}(e), $F(y^{(j)})=U_j\to\OPT$, since
$E_j\to0$. We claim $\dist(y^{(j)},\mathcal S)\to0$. Otherwise a subsequence
stays at distance at least some $c>0$ from $\mathcal S$ and, $X$ being compact,
has a limit point $\bar v\in X$ with $F(\bar v)=\OPT$ and
$\dist(\bar v,\mathcal S)\ge c$, which is impossible. So for all large $j$,
$\dist(y^{(j)},\mathcal S)\le\tau_{\mathrm{TU}}/(2\sqrt{n_c})$ and
$E_j<1/\Omega_{\mathrm{TU}}^2$. Then Lemma~\ref{lem:tu-snap} shows that
step~(ii) finds $\hat x$ with $(\hat x,z)\in\mathcal S$; its value is
$\OPT=a/W$ with $W\le\Omega_{\mathrm{TU}}$, and
$F(\hat x,z)-\beta_j=\OPT-\beta_j\le U_j-\beta_j=E_j<1/\Omega_{\mathrm{TU}}^2
\le1/(\Omega_{\mathrm{TU}}W)$, so the point is returned. This is the argument
of Theorem~\ref{thm:transfer}(b).

(c) By \eqref{eq:tu-setgrowth} and Proposition~\ref{prop:tu-sound}(e),
$g_S\dist(y^{(j)},\mathcal S)^2\le U_j-\OPT\le E_j$, so the two conditions used
in (b) hold as soon as $E_j\le g_S\tau_{\mathrm{TU}}^2/(4n_c)$ and
$E_j\le1/(2\Omega_{\mathrm{TU}}^2)$. Since $g_S\ge\bar L/\kappa_c$ and
$E_j=n_c\bar Lh_j^2/8$, it suffices that
$h_j\le\tau_{\mathrm{TU}}/(n_c\sqrt{\kappa_c})$ and
$h_j\le2/(\Omega_{\mathrm{TU}}\sqrt{n_c\bar L})$, which hold for
$j\ge J_{\rm ex}$ because $h_j=\eta\,2^{-j}$. The quantities
$\log_2(1/\tau_{\mathrm{TU}})$, $\log_2\Omega_{\mathrm{TU}}$, $\log_2\eta$,
$\log_2n_c$ and $\log_2\bar L$ are polynomial in $I$. The last statement
follows from Theorems~\ref{thm:tu-states} and~\ref{thm:tu-approx} and the
remark after Algorithm~\ref{alg:tuexact}.
\end{proof}

\begin{remark}[Unique minimizers without linear programming]\label{rem:tu-cf}
Assume \eqref{eq:tu-setgrowth} with $\mathcal S=\{v^*\}$. Then the hull
variant can return $v^*$ by continued fractions alone. By
Corollary~\ref{cor:tu-height}(a), $v^*_x\in\rho^{-1}\Z^{n_c}$ with
$\rho\le\Delta_\eta R_{\mathrm{TU}}$. By Proposition~\ref{prop:tu-sound}(d)
and Theorem~\ref{thm:tu-states}(c), the retained hull in coordinate $i$
contains $v^*_i$ and has length at most $2a_j+2h_j$, and every retained value
of a discrete coordinate $k$ lies within $a_j$ of $v^*_k$. Two distinct rationals with denominators at most
$\Delta_\eta R_{\mathrm{TU}}$ differ by at least $(\Delta_\eta R_{\mathrm{TU}})^{-2}$,
so once $2a_j+2h_j<(\Delta_\eta R_{\mathrm{TU}})^{-2}$ and $a_j<1$, the rational
of least denominator in each continuous hull is $v^*_i$, and each value set
is $\{v^*_k\}$, since values are integers. Once also
$E_j<\Omega_{\mathrm{TU}}^{-2}$, the candidate passes the test of
step~(iii) of TU-EXACT. A candidate is returned only if it is feasible and
passes this test, so the test, not the reconstruction, carries the proof.
\end{remark}

\subsection{Details for Section~\ref{sec:tu-limits}}\label{app:tu-limits}

\begin{lemma}[TU-GRID keeps order $\sqrt{n_c}$ nodes per coordinate]\label{lem:tu-uniform}
Let $n_c\ge4$ be even, $1\le p\le n_c$ and $\bar L>0$ rational, and consider
the model of Definition~\ref{def:tu-model} with $n_d=0$,
$X=\{x\in[-1,1]^{n_c}:x_1+\dots+x_p\le p\}$, $F(x)=\frac{\bar L}2\norm x^2$,
$\eta=1$ and $\mathcal W=\R^{n_c}$. The row is redundant, so $X$ is the box;
$F$ has the unique minimizer $0$, \eqref{eq:tu-setgrowth} holds with
$\mathcal S=\{0\}$ and $g_S=\bar L/2$, so $\kappa_c=2$, and every admissible
decomposition has a bag containing $\{1,\dots,p\}$. Let
$\nu=\lfloor\sqrt{(n_c-2)/8}\rfloor$. At every level $j$ with
$(\nu+1)h_j\le1$ at which TU-GRID (Algorithm~\ref{alg:tugrid}) filters, both
variants retain $\mathcal D^{(j+1)}_i\supseteq[-(\nu+1)h_j,(\nu+1)h_j]$ for
every $i$, so the grid of level $j+1$ has at least $4\nu+5\ge\sqrt{2(n_c-2)}+1$ nodes in
every coordinate, and the bag containing $\{1,\dots,p\}$ has at least
$(\sqrt{2(n_c-2)}+1)^p\ge n_c^{p/2}$ table entries. For every
$\varepsilon\le\bar L/64$, TU-GRID$(\varepsilon)$ computes the tables of such
a level $j+1$.
\end{lemma}

\begin{proof}
This is the instance of Proposition~\ref{prop:sharp} with $n=n_c$,
$\Lambda=\bar L$ and $g=\bar L/2$, for which the coupling terms vanish and
$\kappa=2$, together with a redundant row. For $\rho>0$ write
$C_\rho=[-\rho,\rho]^{n_c}$.

\emph{Comparison of min-marginals.} Suppose that a level-$j$ domain
$\mathcal D^{(j)}$ contains $C_\rho$ with $\rho\in h_j\Z$ and $\rho\ge h_j$. The
points of $h_j\Z^{n_c}\cap C_\rho$ are feasible grid points, and $0$ is among
them, so $\mu_j=U_j=0$. The uniform grid of mesh $h_j$ on $C_\rho$ is a
grid for the subbox $C_\rho$ in the sense of Section~\ref{sec:grids}, with
corrections $\bar Lw_i^2/8\le\bar Lh_j^2/8$: every coordinate grid contains
$0$ with both adjacent intervals of length $h_j$, and the corrections of a
grid point sum to at most $E_j$. For a node $a$ of coordinate $i$ in
$C_\rho$, the min-marginal $m_i(a)$ of TU-GRID is therefore at most the
min-marginal of this grid at $a$: TU-GRID minimizes over a larger set of
grid points and subtracts $E_j$, which is at least the correction subtracted
at each point.
By Proposition~\ref{prop:sharp} with $U=0$, the latter is at most $0=U_j$
whenever $\lvert a\rvert\le\min\{\rho,h_j\sqrt{(n_c-2)/8}\}$. Hence every cell
of $\mathcal D^{(j)}_i$ with such an endpoint is retained, in both variants.

\emph{Induction, as in the proof of Corollary~\ref{cor:uniformgrid}.} Here
$h_j=2^{-j}$. Put $\rho_j=\min\{1,2(\nu+1)h_j\}$, which lies in $h_j\Z$ and is at
least $h_j$. We show $\mathcal D^{(j)}\supseteq C_{\rho_j}$. For $j=0$,
$\mathcal D^{(0)}=[-1,1]^{n_c}=C_1$ and $\rho_0=1$. Suppose
$\mathcal D^{(j)}\supseteq C_{\rho_j}$. If $(\nu+1)h_j\le1$, then
$(\nu+1)h_j\le\rho_j$, and the cells $[kh_j,(k+1)h_j]$ with $-\nu-1\le k\le \nu$ lie in
$C_{\rho_j}$ and have an endpoint of absolute value at most
$\nu h_j\le\min\{\rho_j,h_j\sqrt{(n_c-2)/8}\}$; they are retained, and
$\mathcal D^{(j+1)}\supseteq C_{(\nu+1)h_j}=C_{\rho_{j+1}}$. If $(\nu+1)h_j>1$, then
$\rho_j=1$, and every cell in $[-1,1]$ has an endpoint of absolute value at
most $1-h_j<\nu h_j$; all are retained, and
$\mathcal D^{(j+1)}\supseteq C_1\supseteq C_{\rho_{j+1}}$. This completes the
induction.

\emph{Counting.} When $(\nu+1)h_j\le1$, the interval $[-(\nu+1)h_j,(\nu+1)h_j]$
contains $4(\nu+1)+1$ points of $h_{j+1}\Z$, and $\nu+1\ge\sqrt{(n_c-2)/8}$ gives
$4\nu+5\ge\sqrt{2(n_c-2)}+1$. The table of the bag containing $\{1,\dots,p\}$ is
indexed by the product of the grids of its coordinates, and
$(\sqrt{2(n_c-2)}+1)^2=2n_c-3+2\sqrt{2(n_c-2)}\ge n_c$. Finally, if
$\varepsilon\le\bar L/64$, the last level $J$ satisfies
$n_c\bar L4^{-J}/8=E_J\le\varepsilon$, so $4^J\ge8n_c$, $J\ge1$ and
$h_{J-1}=2^{1-J}\le(2n_c)^{-1/2}$. Since
$\nu+1\le\sqrt{n_c/8}+1\le\sqrt{2n_c}$ for $n_c\ge1$, the level $J-1$ satisfies
$(\nu+1)h_{J-1}\le1$, and the table of level $J$ is as stated.
\end{proof}

Theorem~\ref{thm:tu-states} gives at most $5+\lfloor2\sqrt{2n_c}\rfloor$ nodes
per coordinate on this instance, so the order $\sqrt{n_c}$ is attained.

\paragraph{The misaligned graded grids after Proposition~\ref{prop:tu-misaligned}.}
With $h=\frac1{16}$ and $\theta=\frac14$, the graded grid of
Definition~\ref{def:graded} of $[0,1]$ with center $\frac3{10}$ has the nodes
\[
0,\ \tfrac{79}{1280},\ \tfrac{51}{320},\ \tfrac{19}{80},\ \tfrac3{10},\
\tfrac{29}{80},\ \tfrac{141}{320},\ \tfrac{689}{1280},\ \tfrac{3381}{5120},\
\tfrac{16649}{20480},\ 1,
\]
and the grid with center $\frac7{10}$ is its reflection $t\mapsto1-t$. The
two grids have only $0$ and $1$ in common. With $m=\frac12$, the left side of
\eqref{eq:tu-gap} is $\frac L4$ at both common nodes, while every interval
of either grid has length at most $\frac{3831}{20480}<\frac15$, so
$d_1+d_2<2\cdot\frac L8\cdot\frac1{25}=\frac L{100}$ there. The smallest positive difference
$y_2-y_1$ is $\delta=\frac{27}{1280}$.

\paragraph{The feasible grid points of Example~\ref{ex:tu-sum}.}
A point of $G^3$ with $x_1+x_2=x_3$ has $x_3\in\{0,1,\frac94,3\}$; since
$2$, $\frac54$, $\frac34$ and $\frac{13}4$ are not in $G$, the feasible grid
points are $(0,0,0)$, $(1,0,1)$, $(0,1,1)$, $(\frac94,0,\frac94)$,
$(0,\frac94,\frac94)$, $(3,0,3)$ and $(0,3,3)$. Their values of
$F-d(x_1)-d(x_2)-d(x_3)$ are $\frac{21}8L$, $\frac{31}{64}L$,
$\frac{31}{64}L$, $\frac{51}{64}L$, $\frac{51}{64}L$, $\frac{175}{64}L$ and
$\frac{175}{64}L$, all positive. The allowance $n_c\bar Lw^2/8$ of
Lemma~\ref{lem:tu-allow}, with $\bar L=L$ and the largest interval
$w=\frac54$ in place of $h_j$, equals $\frac{75}{128}L$ and also gives an
invalid bound: the smallest value of $F$ at these points is $L$, and
$L-\frac{75}{128}L=\frac{53}{128}L>\OPT$.
